\begin{longtable}{|p{1.5cm}|p{5cm}|p{7cm}|}
\caption{Deskripsi \textit{Use Case} Aplikasi \textit{Tracker} ITB Ultra-Marathon}
\label{tab:usecase-deskripsi} \\
\hline
\textbf{ID} & \textbf{Nama \textit{Use Case}} & \textbf{Deskripsi} \\ \hline
\endfirsthead

\caption[]{Deskripsi \textit{Use Case} Aplikasi \textit{Tracker} ITB Ultra-Marathon (lanjutan)} \\
\hline
\textbf{ID} & \textbf{Nama \textit{Use Case}} & \textbf{Deskripsi} \\ \hline
\endhead

\hline
\multicolumn{3}{r}{\textit{Bersambung ke halaman berikutnya}} \\
\endfoot

\hline
\endlastfoot

UC-01 & Melakukan Registrasi Akun &
Aktor dapat mendaftarkan diri ke dalam sistem dengan mengisi data identitas
dan informasi tim. \textit{Use case} ini dapat diakses oleh Peserta dan
Ketua Tim. \\ \hline

UC-02 & Melakukan \textit{Login} Akun &
Aktor melakukan autentikasi untuk mengakses fitur sistem sesuai peran
masing-masing. Dapat dilakukan oleh seluruh aktor yang memiliki akun, yaitu
Peserta, Ketua Tim, Supporter, dan Panitia. \\ \hline

UC-03 & Mengundang \textit{Spectator} &
Peserta dapat mengundang pihak luar (penonton) untuk memantau progresnya
secara individual melalui tautan atau kode unik yang digenerate sistem. \\ \hline

UC-04 & Melihat ETA \textit{Information} &
Peserta dapat melihat estimasi waktu kedatangan ke \textit{checkpoint}
berikutnya maupun ke garis \textit{finish} berdasarkan kecepatan dan jarak
tempuh saat ini. \\ \hline

UC-05 & Melihat \textit{Activity Metrics} &
Peserta dapat memantau data performa aktivitas lari secara \textit{real-time},
meliputi kecepatan, jarak tempuh, dan waktu berlari. \\ \hline

UC-06 & Melakukan \textit{Check-in Checkpoint} &
Peserta melakukan konfirmasi kehadiran saat tiba di titik \textit{checkpoint}
yang telah ditentukan sepanjang rute perlombaan. \\ \hline

UC-07 & Melakukan Pergantian \textit{Relay} &
Peserta pada kategori relay melakukan serah terima dengan anggota tim
berikutnya di titik pergantian yang telah ditetapkan. \\ \hline

UC-08 & Melihat Status Marathon &
Aktor dapat melihat status perlombaan secara keseluruhan, termasuk posisi
peserta, klasemen, dan kondisi umum \textit{event} secara \textit{real-time}.
Dapat diakses oleh Peserta, Ketua Tim, dan Supporter. \\ \hline

UC-09 & Mengundang Peserta &
Ketua Tim dapat mengundang anggota untuk bergabung ke dalam tim relay
yang telah dibentuk di dalam sistem. \\ \hline

UC-10 & Mengundang \textit{Supporter} &
Ketua Tim dapat mengundang pihak pendukung untuk terhubung dengan tim
dan memantau perkembangan perlombaan. \\ \hline

UC-11 & Melihat Tim &
Ketua Tim dapat melihat informasi seluruh anggota tim, termasuk status aktif
pelari, posisi, dan progres masing-masing anggota dalam perlombaan. \\ \hline

UC-12 & Mengirim \textit{Alert Notification} &
Sistem secara otomatis mengirimkan notifikasi peringatan kepada pihak terkait
ketika terdapat kondisi tertentu, seperti peserta melewati \textit{checkpoint}
atau terjadi pergantian relay. \textit{Use case} ini bersifat
\textit{include} dari UC-06 dan UC-07. \\ \hline

UC-13 & Menerima \textit{Alert Notification} &
Aktor menerima notifikasi yang dikirimkan sistem terkait perkembangan
perlombaan. Merupakan perluasan (\textit{extend}) dari UC-11 dan dapat
diterima oleh Ketua Tim, Peserta, dan Supporter. \\ \hline

UC-14 & Memantau Peserta (Individual) &
Penonton dapat memantau posisi dan progres peserta tertentu secara individual
melalui tautan yang dikirimkan oleh peserta yang bersangkutan (UC-03). \\ \hline

UC-15 & Mengelola Data Rute, \textit{Checkpoint}, dan \textit{Switch} &
Panitia dapat mengelola seluruh data operasional perlombaan, termasuk
mengunggah data rute GPX, menetapkan titik \textit{checkpoint}, dan mengatur
titik pergantian relay. \\ \hline

\end{longtable}